\chapter{Conclusiones y Trabajo Futuro}
\label{ch:conclusiones}

Este capítulo presenta las conclusiones derivadas de la \textbf{Revisión Crítica
de Diseño (CDR)} del \textbf{CMOS Microscopy Payload (CMOSM)} del proyecto
INTISAT, así como una identificación clara de las líneas de trabajo futuro
necesarias para la maduración del sistema.

% -------------------------------------------------------------------
\section{Conclusiones generales}

El CMOS Microscopy Payload ha sido diseñado, implementado y validado como un
\textbf{demostrador tecnológico funcional} de microscopía sin lentes en
configuración de contacto, cumpliendo con los objetivos planteados para su fase
actual de desarrollo.

Las actividades de diseño y verificación realizadas permiten concluir que:

\begin{itemize}
    \item El sistema de adquisición basado en un sensor CMOS comercial opera de
    manera estable y reproducible bajo control manual de parámetros críticos.
    \item La microscopía de contacto constituye un enfoque viable para la
    observación de estructuras celulares en muestras biológicas planas, siempre
    que se garantice un adecuado acoplamiento óptico muestra--sensor.
    \item La arquitectura funcional del payload es coherente, modular y alineada
    con las restricciones de sistemas CubeSat académicos.
    \item El registro sistemático de metadatos asegura trazabilidad experimental
    y facilita el análisis posterior de los resultados.
\end{itemize}

% -------------------------------------------------------------------
\section{Evaluación del desempeño del sistema}

Los ensayos experimentales realizados en entorno de laboratorio demuestran que el
CMOSM es capaz de capturar imágenes con suficiente contraste y definición para la
identificación visual de estructuras celulares, tales como paredes celulares en
epidermis vegetal.

Las técnicas de procesamiento computacional evaluadas presentan comportamientos
diferenciados:

\begin{itemize}
    \item La captura directa en modo nominal constituye el mecanismo principal de
    formación de imagen y proporciona resultados consistentes.
    \item La Fourier Ptychographic Microscopy (FPM) ha sido validada únicamente
    como técnica experimental, mostrando mejoras limitadas bajo las condiciones
    actuales de implementación.
    \item La super--resolución basada en aprendizaje profundo mejora la
    legibilidad visual de estructuras ya presentes en la imagen, sin introducir
    información física adicional.
\end{itemize}

Esta evaluación respalda la decisión de diseño de priorizar la simplicidad
óptica y el control físico del sistema sobre técnicas de reconstrucción
computacional complejas.

% -------------------------------------------------------------------
\section{Cumplimiento de objetivos del CDR}

La presente CDR cumple con los objetivos establecidos al:

\begin{itemize}
    \item Presentar un diseño completo y coherente del payload CMOSM.
    \item Justificar las decisiones de diseño adoptadas en los dominios óptico,
    mecánico, electrónico y de software.
    \item Documentar actividades de verificación y validación con resultados
    experimentales reales.
    \item Identificar riesgos técnicos y operacionales de forma explícita.
\end{itemize}

En consecuencia, el sistema se considera \textbf{apto para continuar a fases
posteriores de desarrollo}, dentro del alcance definido para el proyecto
INTISAT.

% -------------------------------------------------------------------
\section{Trabajo futuro}

A partir de los resultados obtenidos y de los riesgos identificados, se proponen
las siguientes líneas de trabajo futuro:

\begin{itemize}
    \item Diseño e implementación de una interfaz muestra--sensor optimizada,
    empleando materiales adecuados para acoplamiento óptico controlado.
    \item Calibración dimensional mediante patrones de referencia, como
    microesferas de tamaño conocido.
    \item Evaluación de estabilidad mecánica y térmica bajo condiciones más
    exigentes.
    \item Integración más profunda con el sistema de OBDH del CubeSat.
    \item Planificación y ejecución de ensayos ambientales (vibración, térmico,
    vacío) en fases posteriores.
\end{itemize}

Estas actividades permitirán incrementar el nivel de madurez tecnológica del
payload y ampliar su aplicabilidad en futuras misiones educativas o científicas.

% -------------------------------------------------------------------
\section{Cierre}

El CMOS Microscopy Payload representa una contribución sólida al desarrollo de
instrumentación científica compacta basada en sensores comerciales y
microscopía computacional. La experiencia adquirida durante esta fase establece
una base técnica y metodológica robusta para la evolución futura del proyecto
INTISAT.
